\chapter{Numbers Eleven to Twenty}
\label{ch:numbers-11-20}

Why are English \emph{eleven} and \emph{twelve} so strange --- why not
``oneteen, twoteen''? German has the identical oddity, and because the two
languages are Germanic twins, learning one \emph{explains} the other. Then, at
thirteen, both settle into a rule that never breaks.
\emph{Practice lessons: \texttt{lessons/GE-C12-*.md}.}

\section{\emph{elf} und \emph{zw\"{o}lf} --- ``one left over,'' ``two left over''}
\label{lesson:elf-zwoelf}

\begin{sounds}
\emph{elf} is said as it looks. \emph{zw\"{o}lf} begins \emph{tsv-} (German
\emph{z} $=$ \emph{ts}, \emph{w} $=$ \emph{v}) and carries the rounded
\textbf{\"{o}}: shape your lips for \emph{oh} and say \emph{eh}.
\end{sounds}

\begin{morphologybox}
\textbf{elf} (11) $\leftarrow$ old Germanic \emph{ainlif} $=$ \emph{ain}
(``one'') $+$ \emph{lif}.\\[3pt]
\textbf{zw\"{o}lf} (12) $\leftarrow$ \emph{twalif} $=$ \emph{twa} (``two'')
$+$ \emph{lif}.\\[3pt]
And \textbf{-lif} means ``\textbf{to leave, to remain}.'' So \emph{elf} $=$
``\textbf{one left over}'' and \emph{zw\"{o}lf} $=$ ``\textbf{two left over}.''
\end{morphologybox}

Left over from \textbf{what}? From \textbf{ten} --- from your ten fingers. You
count to ten, you run out of fingers, and you have \textbf{one left}. Eleven is
literally the leftover.

\begin{cousinweb}
\textbf{English says exactly the same thing}: \emph{eleven} $\leftarrow$
\emph{ain-lif}, \emph{twelve} $\leftarrow$ \emph{twa-lif}. Not similar words
--- \emph{the same words}, inherited side by side. The Germanic-twin pattern
that turned up in \emph{Vater}/\emph{father} and \emph{Wasser}/\emph{water}
reaches even the counting.
\end{cousinweb}

\begin{grammarlens}
\textbf{And then it stops.} The leftover-trick covers only \textbf{two}
numbers. From \textbf{13} on, German turns completely regular --- and so does
English.
\end{grammarlens}

\section{\emph{dreizehn} $\to$ \emph{zwanzig} --- the pattern that never breaks}
\label{lesson:zahlen-13-20}

After the two leftovers, German settles into a rule so regular you can generate
every number without being taught it: \textbf{digit $+$ zehn}. And English does
the identical thing --- because \emph{-teen} \textbf{is} \emph{ten}.

\begin{center}
\begin{tabular}{@{}llll@{}}
\toprule
& German & English & build \\
\midrule
13 & \emph{dreizehn} & thirteen & \emph{drei} $+$ \emph{zehn} \\
14 & \emph{vierzehn} & fourteen & \emph{vier} $+$ \emph{zehn} \\
15 & \emph{f\"{u}nfzehn} & fifteen & \emph{f\"{u}nf} $+$ \emph{zehn} \\
16 & \emph{sechzehn} & sixteen & \emph{sechs} $+$ \emph{zehn} (the \emph{s} drops) \\
17 & \emph{siebzehn} & seventeen & \emph{sieben} $+$ \emph{zehn} (shortened) \\
18 & \emph{achtzehn} & eighteen & \emph{acht} $+$ \emph{zehn} \\
19 & \emph{neunzehn} & nineteen & \emph{neun} $+$ \emph{zehn} \\
20 & \emph{zwanzig} & twenty & $\leftarrow$ \emph{twaintig}, ``two tens'' \\
\bottomrule
\end{tabular}
\end{center}

\begin{sounds}
The \emph{z} of \emph{zehn} and \emph{zwanzig} is \emph{ts}: \emph{TSAYN},
\emph{TSVAHN-tsikh}. The final \emph{-ig} of \emph{zwanzig} is normally said
like the soft \emph{ich}-sound, \emph{-ikh}.
\end{sounds}

Every one of them is a digit $+$ \emph{zehn} --- and the English column is the
same compound with \emph{-teen}, which is just \textbf{ten} again
(\emph{thir-teen} $=$ ``three-ten''). Only two small wearings-down:
\emph{sech(s)zehn} and \emph{sieb(en)zehn} clip a sound for easier speech,
exactly as English \emph{thir-} and \emph{fif-} did to \emph{three} and
\emph{five}.

\begin{morphologybox}
\textbf{zwanzig} (20) $\leftarrow$ \emph{twaintig}, ``\textbf{two tens}'' ---
the \emph{-zig} is another worn-down form of \emph{ten}, matching English
\emph{twen-ty}, whose \emph{-ty} is \emph{ten} as well.
\end{morphologybox}

\begin{culture}
\textbf{The contrast worth noticing.} The Romance sisters all \textbf{break}
their pattern partway through the teens --- Portuguese at 16, French and
Italian at 17, swapping over to a ``ten-and-six'' style.
\textbf{German never breaks.} Two leftovers, then one clean rule all the way to
twenty --- and English marches beside it the entire distance.
\end{culture}
